\chapter{Results}
\label{ch:results}

Every table and figure in this chapter is generated from
\texttt{experiments/results/summary.csv} by \texttt{scripts/gen\_report\_assets.py},
and so is every measured figure quoted in the prose around them: each one is a
command the generator writes into \texttt{tables/numbers.tex} and this chapter
asks for by name, so a re measurement rewrites the sentences as well as the
tables. Nothing here is typed by hand, and a quantity that has not been measured
reads \texttt{pending} rather than being estimated.

Release 1.0.0 opened this chapter with the same claim above a table that was
typed by hand, and whose figures reconciled with no generated artifact anywhere
in the repository. That table is gone, the mechanism that let it exist is gone
with it, and the claim is now checked: a number the generator cannot derive is a
build failure rather than a sentence that quietly stops being true.

The \pnlConfigurations{} rows behind this chapter come from one sweep at one
commit. Timings from release 1.0.0 are not comparable with these: the compiler,
the Jacobi algorithm and the timed region all changed, and the work unit changed
with them. The values did not. Release 1.1.0 keeps
\texttt{PNL\_REDUCTION\_ACCUMULATORS} at one, so every iterate and every
residual here is bit identical to the one release 1.0.0 computed, and the before
and after for each headline number is recorded phase by phase in
\texttt{PROGRESS.md} rather than copied into this page, where it would stop being
true at the next sweep.

\section{The machine}

\input{tables/bandwidth}

Both figures are measured on this hardware with the same STREAM triad kernel,
and both are what the efficiency numbers later in this chapter divide by. The
host figure is taken across all worker threads, since a single core cannot
saturate the memory system of a twenty core part.

\section{Iterations to tolerance}
\label{sec:convergence-results}

These are the hardware free numbers, and they come first for that reason. They
are identical on any machine that runs them.

\input{tables/convergence}

\begin{figure}[htbp]
  \centering
  \includegraphics[width=0.86\textwidth]{iteration_counts.pdf}
  \caption{Iterations to a relative residual of $10^{-8}$ on the largest grid
    where every method converges in reasonable time. The horizontal scale is
    logarithmic: the span from Richardson to conjugate gradient is more than two
    orders of magnitude, and the ordering of the methods is the entire argument
    of Chapter \ref{ch:background} made visible.}
  \label{fig:iteration-counts}
\end{figure}

\begin{figure}[htbp]
  \centering
  \includegraphics[width=0.82\textwidth]{convergence_growth.pdf}
  \caption{Iteration count against grid size on logarithmic axes, so a power law
    appears as a straight line and its exponent as the slope. The $O(n^2)$
    methods and the $O(n)$ methods separate visibly, which is the practical
    content of Young's optimal relaxation result and of the conjugate gradient
    error bound \eqref{eq:cg-bound}.}
  \label{fig:convergence-growth}
\end{figure}

\subsection{Agreement with theory}

The convergence suite checks each of the following as a measurement against a
closed form prediction rather than accepting it as prose. All pass.

\begin{itemize}
  \item \textbf{The Jacobi spectral radius.} The measured asymptotic contraction
        factor agrees with $\cos(\pi h)$ from \eqref{eq:rho-jacobi} to within
        $10^{-3}$ relative at both 31 and 63 grid points per side.
  \item \textbf{The factor of two.} The convergence suite asserts that the ratio
        of Jacobi to Gauss Seidel asymptotic contraction factors is two within
        two percent, which is Young's $\rho_{GS} = \rho_J^2$ from
        \eqref{eq:rho-gs}, and it passes.
  \item \textbf{Forward against backward.} The two sweep directions agree on
        iteration count to within one percent, as similarity of their iteration
        matrices on a symmetric operator requires.
  \item \textbf{Young's optimum is a minimum.} Perturbing $\omega$ by $\pm 0.02$
        and $\pm 0.05$ about the closed form \eqref{eq:omega-opt} costs
        iterations in both directions.
  \item \textbf{The change of order.} The suite asserts that doubling the grid
        multiplies the Jacobi count by more than three, consistent with
        $O(n^2)$, and the optimally relaxed SOR count by less than 2.6,
        consistent with $O(n)$. Both hold.
  \item \textbf{Conjugate gradient is $O(n)$.} In the swept counts the
        conjugate gradient iteration count grows from \pnlCgIterationsSmall{} to
        \pnlCgIterationsLarge{} as the grid doubles from \pnlCgGridSmall{} to
        \pnlCgGridLarge{}, a factor of \pnlCgGrowthFactor{}.
  \item \textbf{Line beats point.} Solving each grid line exactly rather than
        each point reduces the count, and line Jacobi stands to line Gauss
        Seidel in the same ratio of two as the point methods do.
  \item \textbf{The red black penalty.} Red black Gauss Seidel takes
        \pnlRedBlackIterations{} iterations against natural ordering's
        \pnlNaturalIterations{} on a \pnlRedBlackGrid{} grid, a penalty of
        \pnlRedBlackPenalty{} percent. This is the price of the parallelism, and
        it is charged to both sides of the device comparison.
\end{itemize}

The discretisation itself is second order: the suite fits the order across three
grid sizes and asserts it is two, and asserts the error constant against the
$\pi^2/12 = 0.8225$ predicted by the leading truncation term. Both pass. The
fitted values themselves are a property of the convergence run rather than of
the sweep, so they are reported by that suite and not quoted here.

\section{Backend cost on identical work}

Because every backend computes bit identical values, the differences below are
attributable to the execution model and to nothing else.

\input{tables/backend_cost}

\input{tables/backend_cost_verdicts}

Release 1.0.0 read that table as "the three thread models land within 6 percent
of one another", which states the finding backwards. The data cannot separate
them: on the shared work the differences between OpenMP, POSIX threads and
\texttt{std::jthread} are frequently smaller than the run to run spread of the
rows they are computed from, and where they are, the verdicts above say so
instead of printing a percentage. The conclusion that survives is the weaker and
more useful one, that at this granularity the choice between the three is a
choice about ergonomics rather than about performance, and it is supported by an
inability to tell them apart rather than by a measured gap.

\begin{figure}[htbp]
  \centering
  \includegraphics[width=0.86\textwidth]{backend_cost.pdf}
  \caption{Time for a fixed sweep count at the largest size in the backend cost
    block, all backends at twenty workers. A single hue is used because this is
    a comparison of magnitudes along one categorical axis: the bar length
    carries the measurement and the label carries the identity.}
  \label{fig:backend-cost}
\end{figure}

\section{Scaling and the core knee}
\label{sec:knee}

\begin{figure}[htbp]
  \centering
  \includegraphics[width=0.86\textwidth]{scaling_speedup.pdf}
  \caption{Speedup against worker count for a Jacobi sweep on a 1023 by 1023
    grid. The muted diagonal is ideal scaling and is reference chrome rather
    than a measured series. Line style and marker shape repeat the distinction
    that colour makes, so the figure survives greyscale printing.}
  \label{fig:scaling}
\end{figure}

\input{tables/knee}

\subsection{What could and could not be established}

Objective 3 asks for the performance core versus efficiency core knee on this
heterogeneous CPU. The honest position is as follows.

The per processor probe \textbf{could not} classify logical processors into
performance and efficiency groups. Inside WSL2 the guest reports fourteen
uniform cores, identical cache sizes and no hybrid flag, and thread affinity
binds to a virtual processor the hypervisor may place anywhere. The probe times
an identical kernel on each logical processor and reports a split only when the
throughputs separate by a margin dominating the within group spread; the verdict
it recorded for this session is stored in the session manifest.

The knee \textbf{could} be established from the aggregate scaling curve, which
depends only on how many cores are engaged and not on which, and therefore
survives virtualisation. It is located as the two segment least squares split
that minimises total residual.

\subsection{The knee is not where the objective expected it}

The fitted knee sits at \pnlKneeRange{} workers, not at the eight the objective
anticipated, and the backends do not all put it in the same place. Beyond it the
fitted slope is slightly negative: adding workers does not merely stop helping,
it costs a little.

\input{tables/knee_verdicts}

The disagreement between the backends is a result rather than a defect of the
fit, and it is reported as one. Each position carries the bootstrap interval of
the refits behind it, the union of those intervals spans \pnlKneeInterval{}
workers, and a claim about where one backend's knee sits relative to another's
is only worth making where their intervals do not overlap.

The bandwidth probe explains the shape independently. Sweeping the STREAM triad
across worker counts on an idle machine gives

\input{tables/bandwidth_scaling}

The memory system saturates by \pnlTriadPeakWorkers{} workers and falls away
thereafter, and it is at its worst at \pnlTriadFloorWorkers{}. The curve is flat
enough across the middle of its range that neighbouring points are frequently
not separable at the precision of the repetitions behind them, which the spread
column of that table says row by row. A Jacobi sweep, whose access pattern is
less regular than a triad, saturates at the same point or slightly sooner, which
is the range the knee table gives.

The conclusion is that for this workload \textbf{the binding constraint is the
memory system, not the core types}. The heterogeneous core count that Objective 3
anticipated would set the knee never becomes the limiting factor, because the
machine runs out of memory bandwidth long before it runs out of fast cores. That
is a more useful finding than the one the objective was written to look for, and
it would have been invisible to a study that measured speedup without measuring
the memory system separately.

\paragraph{A claim withdrawn.} An earlier run of this probe peaked at exactly
eight workers, and eight is the performance core count of this processor, which
looked like indirect evidence of the topology the guest hides. Repeating the
measurement on an idle machine moved the peak, and the curve is flat enough
across the middle of its range that the location of its maximum is not stable
between runs. The coincidence does not survive and the inference is withdrawn.
It is recorded here rather than deleted because a plausible pattern that fails to
reproduce is worth exactly one sentence of warning to the next person who notices
it.

\section{Pinning}

\input{tables/pinning}

\input{tables/pinning_verdicts}

Under a hypervisor a binding request constrains guest placement and not silicon,
so the expected effect is smaller than the same experiment would show on bare
metal. The measurement is reported for what it is, policy by policy, and most of
those comparisons are not separable at this precision.

\section{The cost of determinism}
\label{sec:reduction-cost}

Section \ref{sec:determinism} argued that reproducible reductions are worth
having. This is what they cost.

\input{tables/reduction_cost}

\input{tables/reduction_cost_verdicts}

Conjugate gradient is used because it performs two global reductions per
iteration and is therefore the method most exposed to the difference. On a
solver with one reduction per sweep the overhead is proportionally smaller.

Read the spread column before the effect column. Release 1.0.0 quoted this block
as a cost of between minus 5 and plus 3 percent and offered it as what
reproducibility is worth, on runs whose own repetitions moved further than that.
The effects here are mostly smaller than the spread of the rows they are
computed from, so what this table establishes is an upper bound rather than a
price: the deterministic reduction does not cost enough on these backends for
this measurement to see it.

\section{Scheduling}
\label{sec:schedule-cost}

\input{tables/schedule_cost}

\input{tables/schedule_cost_verdicts}

This is the measurement behind leaving work stealing out of the jthread pool. On
a uniform stencil sweep a static partition is already balanced, and the extra
bookkeeping of a dynamic schedule buys nothing. The evidence is only as strong as
the verdicts above make it: where a backend's dynamic against static difference
is smaller than the spread of those two rows, that backend contributes nothing to
the argument, and the sentence above rests on the backends where it is larger.

\section{Distributed scaling and the communication fraction}
\label{sec:mpi-results}

\begin{figure}[htbp]
  \centering
  \includegraphics[width=0.86\textwidth]{mpi_scaling.pdf}
  \caption{Speedup against rank count by method. Jacobi exchanges only halo rows
    and scales; conjugate gradient additionally performs two global reductions
    per iteration, each a synchronisation point no amount of local work can
    hide; natural ordering Gauss Seidel is sequentially dependent and is
    included to show what that costs when the semantics are preserved rather
    than quietly replaced.}
  \label{fig:mpi-scaling}
\end{figure}

The MPI backend records halo, reduction, barrier and ordered pass time
separately, so the communication fraction is measured rather than inferred from
the gap between observed and ideal speedup.

\section{The CPU versus GPU comparison}
\label{sec:device-results}

Read this section in the order it is presented. The iteration counts of Section
\ref{sec:convergence-results} come first and are hardware free. The efficiency
column below is the comparable one. The seconds column is a property of this
particular pair of devices.

\input{tables/device_comparison}

\begin{figure}[htbp]
  \centering
  \includegraphics[width=0.82\textwidth]{device_efficiency.pdf}
  \caption{Percentage of each device's own measured STREAM triad bandwidth
    achieved, method by method. This is the fair comparison: it is dimensionless
    and says how well each device is used, rather than which device is larger.
    The outlined bars are the counted byte model against each triad recounted
    the same way, so they sit on the solid bar for a method that is one pass to
    one sweep and above it, as an upper bound, for one that is not.}
  \label{fig:device-efficiency}
\end{figure}

\subsection{Decomposing the speedup}

The wall clock ratio between the two devices is not a single fact. It factors
into two, and the factorisation is the point:

\begin{equation}
  \underbrace{\frac{t_{\text{host}}}{t_{\text{device}}}}_{\text{speedup}}
  \;=\;
  \underbrace{\frac{B_{\text{device}}}{B_{\text{host}}}}_{\text{memory systems}}
  \;\times\;
  \underbrace{\frac{\eta_{\text{device}}}{\eta_{\text{host}}}}_{\text{how well each is used}},
  \label{eq:decomposition}
\end{equation}

with $B$ the measured triad bandwidth of each device and $\eta$ the efficiency
of \eqref{eq:efficiency}. The first factor was bought with the graphics card and
would be the same for any bandwidth bound kernel. The second is the only part
attributable to the implementations.

On the measurements in the table above, the great majority of the observed
speedup is the first factor. That is the honest summary: on a bandwidth bound
stencil sweep this GPU's measured triad is \pnlDeviceOverHostTriad{} times the
host's, and the device is used somewhat more efficiently while doing it. It is a
useful thing to know and considerably less exciting than an undecomposed speedup
figure, which is why it is presented this way.

\paragraph{What each column bounds.} The two bandwidth columns of Table
\ref{tab:device-comparison} are not two estimates of one quantity. Each
percentage divides by the triad counted the same way as its own numerator, so
the read for ownership charge cancels in the ratio and a counted efficiency is
the declared efficiency times the row's pass count over its sweep count. The
byte model therefore moves the achieved bandwidth figure and never the
efficiency. Jacobi is one pass to one sweep, so at \pnlDeviceSize{} unknowns it
reads \pnlJacobiHostEfficiency{} percent on the host and
\pnlJacobiDeviceEfficiency{} percent on the device under both models alike.

For a method with more passes than sweeps the counted model charges every pass a
full stencil pass, and that overstates the traffic: a red black sweep is two
passes that each write half the unknowns, and conjugate gradient's six passes are
one stencil product, two inner products that write nothing and three vector
updates that read what they write. The counted column is an upper bound for those
methods and the declared column a lower one, so the pair brackets the traffic
rather than competing to describe it, and conjugate gradient's counted
\pnlCgHostCounted{} percent on the host and \pnlCgDeviceCounted{} percent on the
device are a bound that is loose rather than a kernel that outran its memory
system. A per pass model that charges each pass what it actually moves is the
honest next step. It wants the traffic model settled first and is release 1.2.0
work, so this release brackets the traffic instead of estimating it.

\input{tables/traffic_model}

\subsection{The crossover, derived}

Because iteration counts and per iteration costs are measured separately, the
point at which one method on one device overtakes another can be computed from
\eqref{eq:crossover} rather than declared. Both sides of that expression are
tabulated in this chapter, so a reader can substitute their own hardware by
replacing the bandwidth term with a triad measurement of their own.
